\documentclass{article}
\usepackage{glossaries}


\newglossaryentry{map}
{
    name = map,
    description = {The playable area in which the player can move around and interact
    with}
}

\newglossaryentry{ScriptOBJ}
{
    name = {scriptable obejct},
    description = {A unity class which stores data}
}

\newglossaryentry{GO}
{
    name = {game object},
    description = {Objects within the game}
}

\newglossaryentry{WS}
{
    name = {world space},
    description = {The in game space}
}

\newglossaryentry{interactable}
{
    name = {interactable},
    description = {An object in the game world which the player can interact with}
}



\begin{document}
\section{General}
In this prototype Communication Design is prioritized over other types of design.
Thus the prototype will focus on communicating the general actions and mechanics 
that the player can do as well as how to do those actions and mechanics.
\section{Mechanics}
What can the player do:
\begin{enumerate}
    \item Move around with WASD.
    \item Sprint using Shift and WASD
    \item Interact with the environment through Space.
    \item Toggling the player's light with R.
\end{enumerate}
Why can the player do these options:
\begin{enumerate}
    \item Explore the map
    \item Run away from enemies or map hazards
    \item To pick up keys, open doors, hide in hiding spots, 
    switching lights on, flipping levers and solving puzzles.
    \item To see places that can only be seen when the light is off
    or to hide from an enemy.
\end{enumerate}
How will these mechanics be balanced
\begin{enumerate}
    \item The player speed will be slower than the base enemy speed
    but still fast enough to not make moving around feel like a drag.
    \item The player sprint speed will be faster than base enemy speed
    but the player will only be able to sprint for a set amount of time.
    \item The player will only be able to interact with one item at a time
    if the player is carrying an item and wants to open a door the player
    will first have to drop the item before opening the door.
    \item The player will not be able to see the map clearly when their
    light is turned off, thus the player will need to toggle between light
    on and light off states in order to use the mechanic.
\end{enumerate}
How will these mechanics be implemented.
\begin{enumerate}
    \item A player script will move the player.
    \item A sprint script will increase the player's speed and keep track of 
    how long the player has been sprinting 
\end{enumerate}

\section{General Description}

\section{Hypothesis}
If important game elements are highlighted
then players will recognize that element faster
because the information hierarchy directs the players attention.
\subsection{Hypothesis in prototype}
In the prototype areas where the player should move are highlighted with
a coloured light. 

Important objects that the player would need to pickup
to continue in the game are also highlighted while the area where they should
be used are highlighted to show the player where they should use the items.

The path that the player should follow can also be seen when the player turns
their light off.
\section{Goals}
Some goals for the prototype
\begin{enumerate}
    \item To create a comprehensive tutorial in which the player learns all 
    the game mechanics and uses them in order to understand the mechanics in 
    such a way that they do not need to constantly be reminded of them.
    \item To create a user interface that can communicate the current state
    the player is in.
    \item To make specific code from the first prototype that can only be used
    in specific situations more general so that the code can be reused easily 
    and is scalable for new features.

\end{enumerate}



\section{Genre}
Exploration. The prototype is of the exploration genre because the player needs
to explore the map they were put in to progress.
Puzzle. The prototype is also of the puzzle genre because the player needs to 
interact with different interactables in order to progress.
\subsection{Game Examples}
\subsubsection*{Portal 1 and Portal 2}
In both Portal 1 and Portal 2 the player is given a portal gun which allows the player
to shoot two portals and walk through them . The player can also pick up items with
the portal gun. In both games you can only have two portals out at a time this leads
the player to think about how they are going to use their portals because of its limited use.
In the prototype only being able to interact with one interactable at a time will have a 
similar effect to the limited use of the portal gun in portal two. This can lead to 
situations where the player has to hide from an enemy but they are currently carrying
an item leading the player to drop the item away from the hiding spot and then hiding
before the enemy catches them.


\section{Communication Design}
The main mechanics of the prototype will be shown to the player through a tutorial.

The moving mechanic will be shown by showing the buttons on the screen while requiring
the player to use them to move through the map.

The sprinting mechanic will be shown by showing the button needed to sprint on screen
and giving the player a long part of the map in which they need to sprint to get through
faster.
The sprinting mechanic cannot be used forever thus a stamina bar will be shown on screen
which will drain to empty if the player sprints for too long. This is to show the player
that they cannot sprint forever.

The interact mechanic will be shown to the player by showing different items they can
interact with while highlighting the important intractables in different ways. The player
will only be able to interact with one interactable at a time.

The game will show the player that they can press R on the keyboard to turn off their light. It will also show
an indication in areas where you need to turn off the light in order to progress. This indication
is a purple diamond in the world space near the hidden objects. When an enemy
is nearby and you need to turn off the light the light will change colour. The colour chosen is
a yellow warning colour for enemies.

Different ways to bring attention to game elements
\begin{itemize}
    \item Shine a light on the important element.
    \item Change the element's colour to make it stand out from the background.
    \item Put instructions in the foreground to bring the element to the player's attention.
    \item 
\end{itemize}

\section{Process}
This prototype will be a continuation from the previous prototype. Thus scripts and assets
from the previous prototype.

The cinemachine package was installed from the unity registry in order to do the camera work. A
cinemachine camera was made to follow the player.

Added the navmeshplus package from https://github.com/h8man/NavMeshPlus in order to do 2d navmesh.
This package helps takes unity navmesh, which is predominantly used for 3d projects, and allows it
to work in a 2d environment.

A change from the first to the second prototype was the picking up of items. In the first prototype
the script was done to only work with items thus the player could only interact with items. In the 
second prototype it was changed to work with any interaction in general by making an interaction 
root script which any other type of interaction inherits from.

In the tutorial the areas between the different mechanic tutorials were extended to allow the users
to read about the mechanics and not just need to experience them.

\section{Reflection}


\end{document}